% ============================================================================
% CAPITULO 4: RESULTADOS Y EVALUACION
% ============================================================================

\section{Resultados y Evaluaci\'on}
\label{sec:resultados}

% TODO: Redactar contenido usando skill tfg-cap4-resultados
% ~15-20 paginas con tablas, graficas y analisis

% ----------------------------------------------------------------------------
% 4.1 Metricas de Evaluacion
% ----------------------------------------------------------------------------
\subsection{M\'etricas de Evaluaci\'on}
\label{subsec:metricas-evaluacion}

% Metricas RAGAs: Faithfulness, Answer Relevancy, Context Precision/Recall
% Metricas personalizadas: Combined Score, Context Overlap, Keyword Coverage
% Confidence Score dinamico (6 factores)
% Justificacion de cada metrica

% ----------------------------------------------------------------------------
% 4.2 Benchmarks RAG
% ----------------------------------------------------------------------------
\subsection{Benchmarks RAG}
\label{subsec:benchmarks-rag}

% --- 4.2.1 Dataset de Evaluacion ---
\subsubsection{Dataset de Evaluaci\'on}
\label{subsubsec:dataset-evaluacion}

% 26 preguntas benchmark (evaluation_dataset.json)
% 115 preguntas test completo (test_queries.txt)
% Categorias: FAQ, contextuales, multi-turn, edge cases

% --- 4.2.2 Resultados Benchmark 115 Preguntas ---
\subsubsection{Resultados del Benchmark de 115 Preguntas}
\label{subsubsec:benchmark-115}

% Tabla: Total 115, calidad 79/84 (94%), avg confidence 0.687
% Distribucion de confidence (histograma)
% Tiempo de respuesta promedio: 3.24s
% Analisis de fallos (5 preguntas fallidas)

% ----------------------------------------------------------------------------
% 4.3 Evolucion de Rendimiento
% ----------------------------------------------------------------------------
\subsection{Evoluci\'on de Rendimiento}
\label{subsec:evolucion-rendimiento}

% Tabla/grafica: v1.0 (0.770) -> v4.1.1 (0.94)
% Mejora total: +22.1%
% Hitos clave: v2.0 (+6.5%), v2.1 (+4.3%), v3.0 (+2.0%), v3.1 (+3.6%), v3.2 (+4.1%)
% Analisis de que cambio produjo cada mejora

% ----------------------------------------------------------------------------
% 4.4 Resultados Ensemble
% ----------------------------------------------------------------------------
\subsection{Resultados del Sistema Ensemble}
\label{subsec:resultados-ensemble}

% Tabla comparativa 4 estrategias:
% Consensus 0.903 (26/26), Routing 0.895, Weighted 0.889, Voting 0.872
% gemma2 base: 0.855 (22/26)
% Analisis: cuando ensemble mejora vs cuando no
% Decision pragmatica: gemma2 directo en produccion

% ----------------------------------------------------------------------------
% 4.5 Resultados Context Tracking
% ----------------------------------------------------------------------------
\subsection{Resultados del Context Tracking}
\label{subsec:resultados-context}

% --- 4.5.1 Benchmark de Persistencia ---
\subsubsection{Benchmark de Persistencia de Contexto}
\label{subsubsec:benchmark-contexto}

% 5 conversaciones criticas, 10 preguntas implicitas
% Tasa de exito: 60% global, 100% conversacion critica
% Comparacion antes/despues del ContextTracker

% --- 4.5.2 Anti-Contaminacion ---
\subsubsection{Sistema Anti-Contaminaci\'on (v4.1.1)}
\label{subsubsec:anti-contaminacion}

% Tabla antes/despues:
% Consultas DB en conv activa: 100% -> 0%
% Contaminacion: ~60% -> ~5%
% Latencia: 50ms -> 5ms
% Tasa exito cambio tema: 40% -> 95%

% ----------------------------------------------------------------------------
% 4.6 Resultados de Seguridad
% ----------------------------------------------------------------------------
\subsection{Resultados de Seguridad}
\label{subsec:resultados-seguridad}

% Tests de injection detection
% Tests de sanitizacion XSS
% Cobertura del SDK de seguridad
% Vulnerabilidades identificadas y mitigadas

% ----------------------------------------------------------------------------
% 4.7 Discusion
% ----------------------------------------------------------------------------
\subsection{Discusi\'on}
\label{subsec:discusion}

% Interpretacion global de resultados
% Fortalezas del sistema
% Debilidades identificadas
% Comparacion con trabajos relacionados (seccion 2.7)
% Validez interna y externa de los resultados
